% 原PDF物理页173–178：原书参考文献、索引和英文摘要。
\SourcePage{173}
\chapter*{参考文献}
\addcontentsline{toc}{chapter}{参考文献}
\begin{description}[leftmargin=2.5em,labelwidth=2em]
\item[{[1]}] P. Deligne：《Weil猜想II》（\emph{La conjecture de Weil II}），\emph{Publ. Math. IHÉS} \textbf{52} (1980)，137--252页。
\item[{[2]}] P. Deligne：《Hodge理论III》（\emph{Théorie de Hodge III}），\emph{Publ. Math. IHÉS} \textbf{44} (1974)，5--78页。
\item[{[3]}] P. Deligne：《\term{l-adic}{$\ell$-进}\term{cohomology}{上同调}中的有限性定理》（\emph{Théorèmes de finitude en cohomologie $\ell$-adique}），收入SGA~$4\frac12$，233--251页。
\item[{[4]}] M. Goresky与R. MacPherson：《奇异空间的Poincaré\term{duality}{对偶}》（\emph{La dualité de Poincaré pour les espaces singuliers}），\emph{C. R. Acad. Sci. Paris} \textbf{284}（1977年6月27日），1549--1551页。
\item[{[5]}] M. Goresky与R. MacPherson：《\term{intersection-homology}{交叉同调}理论》（\emph{Intersection homology theory}），\emph{Topology} \textbf{19} (1980)，135--162页。
\item[{[6]}] M. Goresky与R. MacPherson：《\term{intersection-homology}{交叉同调}II》（\emph{Intersection Homology II}），\emph{Inv. Math.} \textbf{72}, 1 (1983)，77--79页。
\item[{[7]}] R. Hartshorne：《留数与\term{duality}{对偶}》（\emph{Residues and duality}），Lecture Notes in Math. \textbf{20} (1966)。
\item[{[8]}] L. Illusie：《余切\term{complex}{复形}与形变》（\emph{Complexe cotangent et déformations}），Lecture Notes in Math. \textbf{239} (1971)及\textbf{283} (1972)。
\item[{[9]}] M. Raynaud：《\term{sheaf}{层}的Euler--Poincaré示性数与阿贝尔簇的\term{cohomology}{上同调}——据Šafarevič、Ogg和Grothendieck的工作》（\emph{Caractéristique d'Euler-Poincaré d'un faisceau et cohomologie des variétés abéliennes [d'après Šafarevič, Ogg et Grothendieck]}），Séminaire Bourbaki 286 (1964/65)，New York: Benjamin, 1966。
\item[{[10]}] J.-L. Verdier：《\term{abelian-category}{阿贝尔范畴}，状态0》（\emph{Catégories abéliennes, état 0}），收入SGA~$4\frac12$，262--308页。

注意，这里的某些符号约定与原始油印讲义中的约定不同。
\end{description}

\SourcePage{174}
\section*{缩写}
\begin{description}[leftmargin=4em,labelwidth=3.5em]
\item[{SGA}] Bois-Marie代数几何讨论班（\emph{Séminaire de géométrie algébrique du Bois-Marie}）；除SGA~2外，均收入Lecture Notes in Mathematics丛书。
\item[{SGA 4}] LN 269（I--IV），LN 270（V--VIII），LN 305（IX--XIX）。
\item[{SGA $4\frac12$}] LN 569。
\item[{SGA 5}] LN 589。
\item[{SGA 7}] LN 288（I--IX），LN 340（X--XXII）。
\item[{EGA}] A. Grothendieck与J. Dieudonné著《代数几何基础》（\emph{Éléments de géométrie algébrique}）。EGA~4第三部分载于\emph{Publ. Math. IHÉS} \textbf{28} (1966)。
\end{description}

\SourcePage{175}
\chapter*{原书术语索引}
\addcontentsline{toc}{chapter}{原书术语索引}
\begin{description}[leftmargin=2em,style=nextline]
\item[{\term{admissible}{容许}}] 态射：1.2.3；\term{abelian-subcategory}{阿贝尔子范畴}：1.2.5；短\term{exact-sequence}{正合列}：1.2.3。
\item[{\term{cohomological-amplitude}{上同调幅度}}] 4.2.1。
\item[{笛卡尔}] 3.2.5，3.2.6。
\item[{范畴}] \term{derived-category}{导出范畴}：1.1.3；\term{filtration}{滤过}、双\term{filtered-derived-category}{滤过导出范畴}：1.1.4例1，3.1.1；\term{exact-category}{正合范畴}：1.1.4；\term{triangulated-category}{三角范畴}：1.1.1。
\item[{\term{heart}{心}（$t$-范畴的）}] 1.3.1。
\item[{\term{cone}{锥}}] 1.1.1。
\item[{\term{constructible}{可构造}}] $\mathcal S$-\term{constructible}{可构造}：2.1.13；$(\mathcal S,\mathcal L)$-\term{constructible}{可构造}：2.2.10。
\item[{九图}] 1.1.11。
\item[{\term{distinguished-triangle}{区别三角}}] 1.1.1。
\item[{\term{extension}{扩张}（在其下稳定）}] 1.2.6。
\item[{\term{filtration}{滤过}}] 愚笨\term{filtration}{滤过}：3.1.7；\term{canonical-filtration}{典范滤过}：3.1.5；\term{weight-filtration}{权滤过}：5.3.5。
\item[{局部$\mathcal C$-性}] 3.2.3。
\item[{\term{mixed}{混合}}] 5.1.5。
\item[{八面体}] 1.1.6。
\item[{\term{geometric-origin}{几何来源}}] 6.2.4。
\end{description}

\SourcePage{176}
\begin{description}[leftmargin=2em,style=nextline]
\item[{\term{perverse-sheaf}{反常层}}] 2.1.5，2.1.13；在2.2.1、2.2.17、4.0.1、4.0.2中推广。
\item[{\term{perversity}{反常度}}] \term{duality}{对偶}\term{perversity}{反常度}：2.1.16；\term{perversity}{反常度}函数：2.1.1，2.2.1。
\item[{\term{weight}{权}}] 逐点\term{weight}{权}：5.1.5；\term{weight}{权}$\le w$、$\ge w$：5.1.8，在5.1.11、5.1.12中推广。
\item[{\term{external-tensor-product}{外张量积}}] 4.2.7。
\item[{延拓}] 1.4.13。
\item[{\term{pure}{纯}}] \term{pointwise-pure}{逐点纯}：5.1.5；\term{weight}{权}为$w$的\term{pure}{纯}对象：5.1.8。
\item[{实现}] 3.1.9。
\item[{\term{recollement}{粘合}（\term{t-structure}{$t$-结构}的）}] 1.4.10。
\item[{\term{stratum}{层片}}] 2.1.1。
\item[{$t$-范畴}] 1.3.1。
\item[{\term{t-exact}{$t$-正合}}] 1.3.16；多元函子的\term{t-exact}{$t$-正合}性：1.3.20。
\item[{\term{t-structure}{$t$-结构}}] 1.3.1；\term{duality}{对偶}\term{t-structure}{$t$-结构}：1.3.2(iii)；\term{induced-t-structure}{诱导$t$-结构}：1.3.19；自然\term{t-structure}{$t$-结构}：1.3.2；\term{perversity}{反常度}为$p$的\term{t-structure}{$t$-结构}：2.1.4；由平移得到的\term{t-structure}{$t$-结构}：1.3.2(ii)。
\item[{\term{shift-functor}{平移函子}}] 1.1.1。
\item[{\term{truncation}{截断}}] 1.3.2(i)，1.3.3。
\end{description}

\SourcePage{177}
\chapter*{原书符号索引}
\addcontentsline{toc}{chapter}{原书符号索引}
\begin{itemize}
\item $\mathcal D$为一个\term{triangulated-category}{三角范畴}，通常带有\term{heart}{心}为$\mathcal C$的\term{t-structure}{$t$-结构}。
\[
\begin{gathered}
\mathcal D^{\ge a},\mathcal D^{\le b}:\ 1.3.1;\qquad
\mathcal D^{[a,b]}:=\mathcal D^{\ge a}\cap\mathcal D^{\le b};\\
\mathcal D^b:\ 1.3.12;\qquad
\mathcal D_F,\mathcal D_F^b:\ 3.1.5.
\end{gathered}
\]
\item $D(\ldots)$为依赖于所列数据的\term{triangulated-category}{三角范畴}。
\[
\begin{gathered}
D(X,\mathcal O):\ 1.4.1;\quad
D_{\mathcal S}:\ 2.1.13;\quad
D_c:\ (2.1.13),\ 2.2.1;\\
D_{\mathcal S,\mathcal L}:\ 2.1.10;\qquad DF,DF^2:\ 3.1.1;\\
D_c^b(X,\Z_\ell):\ (2.2.14.1);\\
D_c^b(X,\Q_\ell),\ D_c^b(X,E_\lambda),\ D_c^b(X,\overline{\Q}_\ell):\ 2.2.17.
\end{gathered}
\]
\item $p$为\term{perversity}{反常度}（2.1.1，2.2.1），$p^*$为其\term{duality}{对偶}\term{perversity}{反常度}；$p_s$：2.2.1。

作为左上标时，$p$表示$t$-范畴的\term{heart}{心}之间的函子：${}^p\mathrm T$，见1.3.17(i)；或者表示相对于某个\term{t-structure}{$t$-结构}的对象，例如${}^p\tau,{}^pH^i$。
\item $\tau$为\term{truncation}{截断}函子。
\[
\tau_{\ge a},\tau_{\le b}:\ 1.3.3;\quad
\tau_{[a,b]}:\ 1.3.5;\quad
\tau^F,\tau^U:\ 1.4.13.
\]
\end{itemize}
\noindent\textbf{其他记号}
\begin{description}[leftmargin=6em,labelwidth=5.5em]
\item[{$\mathrm D$}] Verdier\term{duality}{对偶}。
\item[{$F[n]$}] 平移后的\term{filtration}{滤过}：3.1.8。
\item[{$G$}] $\mathcal D_F^b\to C^b(\mathcal D)$：3.1.7。
\item[{$H^0$}] 1.3.6。
\item[{$\operatorname{Hom}_\alpha$}] 3.2.5。
\item[{$j_{!*}$}] 1.4.14.1，1.4.22，2.1.7。
\item[{$M(X,\mathcal O)$}] 1.4.1。
\item[{$n,n^*,n^+$}] 3.3.1。
\item[{$\mathrm{r\acute eal}$}] 3.1.9。
\item[{$\operatorname{tot} A$}] 3.2.7。
\item[{$\delta(F,G)$}] 3.1.2。
\item[{$\Delta,\Delta^+$}] 3.2.6。
\item[{$\boxtimes$}] 4.2.7。
\end{description}

\SourcePage{178}
\chapter*{原书英文摘要中译}
\addcontentsline{toc}{chapter}{原书英文摘要中译}
这几卷收录了Luminy会议“奇异空间上的分析与拓扑”（\emph{Analyse et topologie sur les espaces singuliers}）的讲义。报告围绕以下主题展开：\term{perverse-sheaf}{反常层}，\term{intersection-homology}{交叉同调}及其应用，\term{intersection-homology}{交叉同调}与\term{vanishing-cycles}{消失轮换}上的Hodge结构，局部与整体Euler--Poincaré示性数，以及Bernstein多项式。
